\section{Introduction}
\label{sec:introduction}
\begin{figure}[t!]
\centering
    % \includegraphics[width=\linewidth]{attachments/figure1.pdf} % 이미지 로드 주석 처리
    % \framebox[.95\linewidth]{\rule{0pt}{10cm}} % 너비는 \linewidth, 높이는 5cm인 빈 박스 (높이는 조절 가능)
    \includegraphics[width=\linewidth]{attachments/figure1-img.pdf}
    \caption{Overview of existing video summarization evaluation approaches and our proposed method, QEVA.
(Top) Traditional reference-based metrics rely solely on text-to-text comparisons between candidate and human-written reference summaries, often failing to capture nuanced semantic and multimodal content. (Bottom) QEVA leverages a fully reference-free multimodal question-answering pipeline (Coverage QA, Factuality QA, Chronology QA) to directly evaluate candidate summaries against source videos, enabling a more comprehensive and semantically grounded assessment.
\textbf{Takeaway:} QEVA provides a more accurate and scalable alternative by eliminating the reliance on human-written reference summaries and directly assessing summaries against source video content.}
    \label{fig:vtsscore_overview}
  \end{figure}

Text-based video summarization has become increasingly crucial due to the explosive growth of video content and significant advances in Video-Large Multimodal Models (Video-LMMs). Despite substantial progress in generating comprehensive textual summaries from videos, insufficient attention has been paid to reliably evaluating the quality of these summaries. Currently, evaluation primarily relies on reference-based metrics such as ROUGE or METEOR, which compare generated summaries to reference texts produced by human annotators. However, acquiring accurate and detailed reference summaries for videos requires considerable human effort and resources, making the process both costly and inefficient at scale. Consequently, the absence of efficient evaluation approaches has become a major obstacle to the advancement of video summarization research.  

\begin{figure*}[t]
    \centering
    \includegraphics[width=\textwidth]{attachments/figure2-img.pdf}
    \caption{Detailed illustration of QEVA's multimodal question-answering methodology. Given a video and candidate summary, QEVA evaluates the summary across three distinct dimensions: (1) Coverage (whether the summary comprehensively covers key video content), (2) Factuality (accuracy of the information presented), and (3) Chronology (i.e.,chronological fidelity; correctness of event ordering). QEVA generates tailored QA pairs through a structured pipeline involving question generation, question answering, and answer correctness checking.
\textbf{Takeaway:} QEVA explicitly decomposes video summary quality into three complementary dimensions, enabling nuanced and interpretable evaluation.}
    \label{fig:qeva_pipeline}
    \end{figure*}

While reference-free metrics have gained considerable attention in text summarization tasks, these methods typically compare the generated summaries directly with the source text, leveraging linguistic similarities or embeddings. Unfortunately, extending this concept to videos is inherently challenging due to the fundamental modality gap - videos are spatio-temporal and multimodal, and thus cannot be directly compared to textual summaries in a straightforward manner. Moreover, recent attempts at multimodal evaluation, such as CLIPScore or LLM-based comparisons used in benchmarks like MLVU, still rely heavily on reference summaries or fail to reliably capture important aspects such as chronological fidelity or factual accuracy. Thus, there is a pressing need for a novel, effective, and fully reference-free evaluation paradigm tailored specifically for video summarization.  

To address these challenges, we propose a \textbf{Q}uestion-answering based \textbf{E}valuation metric for \textbf{V}ideo summ\textbf{A}rization, QEVA. QEVA is built upon the intuitive principle that a high-quality video summary should be able to substitute the original video content effectively. Based on this principle, we identify three critical dimensions for evaluating summaries: Coverage (capturing all essential content), Chronology (chronological fidelity; preserving the order of events), and Factuality (ensuring factual correctness). QEVA leverages Video-LMMs to automatically generate relevant questions from the original video content across these dimensions and employs LLMs to answer these questions using only the generated summaries. By measuring the accuracy of these answers, QEVA quantitatively assesses summary quality without the need for any human reference summaries.

Furthermore, we present MLVU(VS)-Eval, a new annotated evaluation dataset derived from the MLVU benchmark, containing 800 summaries generated by Video-LMMs such as GPT-4o, QwenVL, InternVL, and Video-LLaVA. Each summary has been rigorously annotated by multiple human evaluators according to Coverage, Chronology, and Factuality, demonstrating strong inter-annotator agreement (Krippendorff's $\alpha$ = 0.68). Our comprehensive experiments show that QEVA consistently exhibits the highest correlation with human judgments compared to existing reference-based and multimodal evaluation methods (e.g., ROUGE, METEOR, BERTscore).

By introducing QEVA and the MLVU(VS)-Eval dataset, this work pioneers the first systematic exploration of reference-free evaluation for video-to-text summarization. Our approach not only addresses the pressing scalability and practicality issues but also sets a clear foundation for future advancements in multimodal summarization research, facilitating rapid and reliable assessment of Video-LMMs in real-world applications, including content platforms, news summarization, and automated content generation services.


% 